Pelaksanaan penelitian direncanakan berlangsung secara fleksibel dan menyesuaikan dengan ketersediaan waktu serta kebijakan mitra kerja sama.  
Peneliti mengajukan rentang waktu awal pelaksanaan pada tanggal 2 Februari hingga 8 Februari 2026.  
Penjadwalan detail akan ditentukan melalui koordinasi lanjutan setelah proses persetujuan dan wawancara dengan pihak PUSBISINDO.  

Rencana waktu pelaksanaan secara umum disajikan pada Tabel~\ref{tab:waktu-penelitian}.  

\begin{table}[H]
\centering
\caption{Rencana Waktu Pelaksanaan Penelitian}
\label{tab:waktu-penelitian}
\begin{tabular}{p{6cm} p{6cm}}
\hline
\textbf{Tahapan} & \textbf{Perkiraan Waktu} \\
\hline
Pengajuan dan koordinasi awal & 2--8 Februari 2026 \\
Wawancara dan pembahasan teknis & Menyesuaikan kesepakatan bersama \\
Pengambilan data & Menyesuaikan jadwal mitra dan partisipan \\
Pengolahan dan analisis data & Setelah proses pengambilan data selesai \\
\hline
\end{tabular}
\end{table}

Tempat pelaksanaan pengambilan data direncanakan bersifat fleksibel dan menyesuaikan dengan kenyamanan serta rekomendasi dari pihak PUSBISINDO.  
Beberapa lokasi yang dipertimbangkan oleh peneliti disajikan pada Tabel~\ref{tab:lokasi-penelitian}.  

\begin{table}[H]
\centering
\caption{Alternatif Lokasi Pengambilan Data}
\label{tab:lokasi-penelitian}
\begin{tabular}{p{4cm} p{7cm}}
\hline
\textbf{Lokasi} & \textbf{Keterangan} \\
\hline
MIRE Hub & Jl. Brigadir Jend. Katamso No.19B, Cihaur Geulis, Kec. Cibeunying Kaler, Kota Bandung, Jawa Barat 40122 \\
Area Depok & Lokasi alternatif dengan dukungan komunitas Juru Bahasa Isyarat yang dikenal peneliti \\
Delasic Studio & Jl. Mars Dirgahayu, Cibeunying, Kec. Cimenyan, Kab. Bandung, Jawa Barat 40191 (dalam tahap konfirmasi) \\
Lokasi lain & Coworking space atau lokasi lain yang direkomendasikan oleh PUSBISINDO \\
\hline
\end{tabular}
\end{table}

Peneliti terbuka untuk menggunakan lokasi yang disarankan atau disediakan oleh pihak PUSBISINDO apabila dinilai lebih sesuai dan nyaman bagi partisipan.  